\textbf{What the foundation models do well.} The Chronos models were never shown Indian temperature data, yet Chronos-2 and Chronos-Bolt beat persistence and climatology on CRPS pooled over leads 1-7 (Chronos-T5 did not beat persistence at the focus regions). They also give usable probabilities: the hot-day reliability curves for Chronos-2 and Chronos-2+cov lie close to the diagonal, and their heatwave ROC areas are above 0.96. For an agency office that has to cover many districts with little data engineering, this is the main practical result. A single pretrained model, fed only the recent temperature history, produces calibrated-enough 7-day probabilities anywhere on the grid without a training step.

\textbf{Where they fall short.} Models trained on 64 years of local data remain better. The quantile LSTM and the AR model lead on CRPS, hot-day BSS and interval coverage. The gap is small at lead 1 and grows with lead time. The case studies in Fig.~\ref{fig:cases} show why. When a heatwave was already underway, every model, trained or not, forecast a return towards normal within a few days while the heat actually persisted. The trained baselines learn how fast anomalies at a given place and season decay. A zero-shot model has to infer that from 512 days of context, and it tends to relax too fast and with intervals that are too narrow (coverage \CovChronostwocov{} against the 0.80 target for Chronos-2+cov).

\textbf{Covariates.} Adding Tmin and the known seasonal normal to Chronos-2 improved the 2015-2018 validation CRPS clearly and helped in \CovBetterYears{} of the \CovYearsTotal{} test years. It was slightly worse in 2019 and clearly worse in 2023 (CRPS \CRPSChronostwocovYtwozerotwothree{} against \CRPSChronostwoYtwozerotwothree{}). The 2023 season had the second lowest mean Tmax anomaly of the six test years and several sharp drops in Tmax, and giving the model the normal as a future covariate seems to pull forecasts towards it during such cool spells. Relative humidity from the station proxy gave a further small gain at the focus regions (CRPS \FocusCRPSChronostwocovRH{} against \FocusCRPSChronostwocov{}). These results suggest covariates matter, but one validation period is not enough to choose them. A deployment would need a longer rolling check.

\textbf{Generative sampling.} Chronos-T5, the only model that generates sample paths token by token, was the weakest Chronos variant on continuous scores at the focus regions and by far the slowest (about \SpeedChronosTfive{} forecasts per second on the laptop, against about \SpeedChronostwo{} for Chronos-2). With only \TFiveSamples{} paths its probabilities are coarse. The share of paths above a threshold and the probability read from the fitted quantiles were almost the same (correlation \TFiveCorr{}), so sample paths add nothing that quantiles do not already give for this task. Its heatwave alerts tied with Chronos-Bolt for the best CSI at the focus regions, but that rests on 105 events and should be read with care.

\textbf{Limits of the grid.} The 1\textdegree{} IMD grid averages station values over about 100\,km, which lowers peaks. That is why the Mumbai and Pune cells never met the IMD rule in 2019-2024, while the point series met it on \LocPointHW{} region-days in the same seasons. Station rules applied to gridded data underestimate heatwave frequency. IMD also asks for the criteria to hold at two stations for two consecutive days. A single cell cannot test the first part; the 2-day variant in Table~\ref{tab:events} shows the second. Severe heatwaves were too rare on the grid (\NSevTest{} days in the test period) to verify the Red level at all.

\textbf{Station proxy.} Open-Meteo point data come from ERA5 reanalysis, not from a thermometer at the site. At Mumbai the proxy runs 2.6\,\textdegree C cooler than the grid cell on average, likely because the reanalysis cell at the coast includes sea surface. The point correction is therefore a correction to another model, and the results in Table~\ref{tab:localization} should be repeated once the use case's IoT weather stations are running. The loader and QC code for that are already in place.

\textbf{Humid heat.} The Tmax-only rule flagged \TmaxhwMumbai{} heatwave days at Mumbai in 1991-2024. Over the same seasons the heat index passed the NOAA Danger level on \HumidMumbai{} other days, about \HumidPerYearMumbai{} a year, against \HumidPerYearPune{} a year at inland Pune. This supports IMD's stated plan to fold humidity and warm nights into its criteria \cite{imd2026criteria}. With minimum humidity as the proxy for humidity at the time of Tmax, the heat index never reached IMD's experimental Orange band ($\ge$46\,\textdegree C). The highest value was \HImaxMumbai\,\textdegree C at Mumbai. So these counts depend on which humidity value and which threshold are used.
